\documentclass[11pt]{article}

\usepackage[margin=0.95in]{geometry}
\usepackage{amsmath,amssymb,amsthm,mathtools,booktabs,enumitem,microtype}
\usepackage{xcolor}
\usepackage[colorlinks=true,linkcolor=blue!50!black,urlcolor=blue!50!black,citecolor=blue!50!black]{hyperref}

\theoremstyle{plain}
\newtheorem{theorem}{Theorem}
\newtheorem{lemma}[theorem]{Lemma}
\newtheorem{proposition}[theorem]{Proposition}
\newtheorem{corollary}[theorem]{Corollary}

\theoremstyle{definition}
\newtheorem{definition}[theorem]{Definition}

\theoremstyle{remark}
\newtheorem{remark}{Remark}

\DeclareMathOperator{\tr}{tr}
\DeclareMathOperator{\rank}{rank}
\DeclareMathOperator{\diag}{diag}
\DeclareMathOperator{\GL}{GL}

\title{\bfseries Resolution of the Three Similarity Matching PCA Conjectures}
\author{}
\date{\vspace{-1.6em}%
\small
\begin{tabular}{@{}c@{}}
Target paper: V.~Centorrino, F.~Bullo, and G.~Russo,\\
\emph{Similarity Matching Networks: Hebbian Learning and Convergence Over Multiple Time Scales},\\
\href{https://arxiv.org/abs/2506.06134}{arXiv:2506.06134} (v1, June 2025; \emph{Neural Computation} 38(5):725--764, April 2026)\\[0.35em]
Lean~4 companion formalization: \texttt{lean\_proofs/Conjectures/SimMatchingPCA.lean}\\[0.35em]
November 2026
\end{tabular}}

\begin{document}

\maketitle

\begin{abstract}
We resolve all three open conjectures of Centorrino, Bullo, and Russo (\href{https://arxiv.org/abs/2506.06134}{arXiv:2506.06134}, \emph{Neural Computation} 2026) on the continuous-time gradient flow of the similarity-matching PCA objective $S_3(W) = 2\|W\|_F^2 - 3\tr((WCW^\top)^{2/3})$ for $C \succ 0$.
First, we prove Conjecture~1 (full-rank invariance and forward completeness): along the flow, the singular values of $W(t)$ remain uniformly bounded in $[\min\{\sigma_{\min}(W_0), \lambda_{\min}(C)\},\, \max\{\sigma_{\max}(W_0), \lambda_{\max}(C)\}]$ for all $t \ge 0$.
Second, we prove the intended form of Conjecture~2 (every global minimum has full row rank and spans a top-$m$ eigenspace of $C$, no rank-deficient point is even a local minimum, and every non-global full-rank stationary point is a strict saddle with explicit negative curvature $4(1 - \beta/\alpha) < 0$ and positive radial curvature $8/3$), while noting a matrix-permutation typo in the printed stationary set of arXiv~v1 and attributing related landscape results to Pehlevan--Sengupta--Chklovskii (2018) and Lipshutz--Lipshutz (2026).
Third, we decisively refute Conjecture~3: $H = U^\top \Lambda_1^\top (\Lambda_1 U \Lambda_2 U^\top \Lambda_1^\top)^{-1/3} \Lambda_1 U$ is diagonal if and only if $U^\top (\Lambda_1^\top \Lambda_1) U$ is diagonal, independently of $\Lambda_2$. In particular, whenever $n - m \ge 2$, any orthogonal rotation inside $\ker(\Lambda_1)$ yields a counterexample for every $\Lambda_1, \Lambda_2$. We also observe that the paper's trace-maximization lemma nonetheless survives via Horn's log-majorization.
\end{abstract}

\section{Setup and Overview of Verdicts}

Let $C \in \mathbb{R}^{n \times n}$ be symmetric positive definite ($C \succ 0$) with eigenvalues $\lambda_1 \ge \dots \ge \lambda_n > 0$, let $1 \le m \le n$, and let $\mathcal{R}_{m,n} := \{W \in \mathbb{R}^{m \times n} : \rank(W) = m\}$ (for $m > n$, $\mathcal{R}_{m,n} = \varnothing$).
On $\mathcal{R}_{m,n}$, the similarity-matching cost and its continuous-time gradient flow are
\begin{equation}\label{eq:flow}
S_3(W) = 2\|W\|_F^2 - 3\tr\!\bigl((WCW^\top)^{2/3}\bigr),
\qquad
\dot{W} = -\nabla S_3(W) = -4W + 4(WCW^\top)^{-1/3}WC.
\end{equation}
On all of $\mathbb{R}^{m \times n}$, $S_3(W)$ extends continuously via the positive semidefinite matrix $2/3$-power.

\begin{itemize}[leftmargin=2em,itemsep=2pt]
\item \textbf{Conjecture~1 (\texttt{conj\_full\_row\_rank}): Proved} (Theorem~\ref{thm:c1}).
\item \textbf{Conjecture~2 (\texttt{conj\_full\_strict\_saddle}): Intended statement proved} (Theorem~\ref{thm:c2}); the printed stationary set in arXiv~v1 contains a row-vs.-column permutation typo producing non-stationary points.
\item \textbf{Conjecture~3 (\texttt{conj\_lin\_prog}): Disproved}, with an exact if-and-only-if criterion (Theorem~\ref{thm:c3}).
\end{itemize}

\section{Conjecture 1: Full-Rank Invariance and Forward Completeness}

\begin{theorem}[Uniform singular-value envelope and global existence]\label{thm:c1}
Let $C \in \mathbb{R}^{n \times n}$ be symmetric positive definite with $c := \lambda_{\min}(C) > 0$ and $d := \lambda_{\max}(C) \ge c$, and let $1 \le m \le n$. For every initial condition $W_0 \in \mathcal{R}_{m,n}$, the unique solution $W(t)$ of~\eqref{eq:flow} exists for all $t \ge 0$, satisfies $W(t) \in \mathcal{R}_{m,n}$ for all $t \ge 0$, and obeys
\begin{equation}\label{eq:sv-envelope}
\min\{\sigma_{\min}(W_0),\, c\} \le \sigma_{\min}(W(t)) \le \sigma_{\max}(W(t)) \le \max\{\sigma_{\max}(W_0),\, d\}
\qquad (\forall t \ge 0).
\end{equation}
\end{theorem}

\begin{proof}
Because $A(W) := WCW^\top \succ 0$ on the open set $\mathcal{R}_{m,n}$ and $A \mapsto A^{-1/3}$ is real-analytic on $\mathbb{S}^m_{++}$, the right-hand side of~\eqref{eq:flow} is real-analytic on $\mathcal{R}_{m,n}$. Hence~\eqref{eq:flow} has a unique maximal solution $W : [0, T_{\max}) \to \mathcal{R}_{m,n}$.

Set $G(t) := W(t)W(t)^\top \succ 0$ and $A(t) := W(t)CW(t)^\top \succ 0$. Differentiating $G(t)$ along~\eqref{eq:flow} yields
\begin{equation}\label{eq:gram-ode}
\dot{G} = \dot{W}W^\top + W\dot{W}^\top = -8G + 4A^{-1/3}A + 4AA^{-1/3} = -8G + 8A^{2/3}.
\end{equation}
Let $\ell(t) := \lambda_{\min}(G(t)) > 0$ and $L(t) := \lambda_{\max}(G(t)) > 0$. Since $cI_n \preceq C \preceq dI_n$, for every $x \in \mathbb{R}^m$ we have
\[
c\ell(t)\|x\|^2 \le c x^\top G(t) x \le x^\top A(t) x \le d x^\top G(t) x \le dL(t)\|x\|^2,
\]
i.e., $c\ell(t) I_m \preceq A(t) \preceq dL(t) I_m$. By spectral mapping of the strictly increasing function $s \mapsto s^{2/3}$ on $(0, \infty)$ (requiring no L\"owner--Heinz theorem because the bounds are scalar multiples of $I_m$),
\begin{equation}\label{eq:A-bounds}
(c\ell(t))^{2/3} I_m \preceq A(t)^{2/3} \preceq (dL(t))^{2/3} I_m.
\end{equation}

\emph{Lower bound.} Fix any strict sublevel $b$ with $0 < b < \min\{\ell(0), c^2\}$. Suppose for contradiction that $\{t \in [0, T_{\max}) : \ell(t) \le b\}$ is nonempty, and let $t_*$ be its infimum. By continuity, $t_* > 0$, $\ell(t_*) = b < c^2$, and $\ell(t) > b$ for all $t \in [0, t_*)$. Choose a unit eigenvector $v \in \mathbb{R}^m$ with $G(t_*)v = bv$, and consider the smooth scalar function $\phi(t) := v^\top G(t) v$. Then $\phi(t) \ge \ell(t)$ for all $t \in [0, T_{\max})$, $\phi(t_*) = b$, and by~\eqref{eq:gram-ode}--\eqref{eq:A-bounds},
\[
\phi'(t_*) = -8b + 8v^\top A(t_*)^{2/3} v \ge 8\bigl((cb)^{2/3} - b\bigr) > 0,
\]
since $0 < b < c^2 \iff b^3 < c^2 b^2 = (cb)^2 \iff b < (cb)^{2/3}$. (Equivalently, for $0 < h < 1/8$, $(1 - 8h)G(t) + 8hA(t)^{2/3} \succeq [(1-8h)\ell(t) + 8h(c\ell(t))^{2/3}]I_m$, so Weyl's inequality yields the lower right Dini derivative $D_+ \ell(t) \ge 8((c\ell(t))^{2/3} - \ell(t))$ regardless of eigenvalue multiplicities.) Since $\phi(t_*) = b$ and $\phi'(t_*) > 0$, we have $\ell(t_* - h) \le \phi(t_* - h) < b$ for sufficiently small $h > 0$, contradicting the minimality of $t_*$. Thus $\ell(t) > b$ for all $t \in [0, T_{\max})$. Letting $b \uparrow \min\{\ell(0), c^2\}$ yields $\ell(t) \ge \min\{\ell(0), c^2\}$.

\emph{Upper bound.} Symmetrically, fix any $B > \max\{L(0), d^2\}$ and suppose $t^* := \inf\{t \in [0, T_{\max}) : L(t) \ge B\} < T_{\max}$. Picking a unit eigenvector $v$ with $G(t^*)v = Bv$ and setting $\psi(t) := v^\top G(t) v \le L(t)$ gives $\psi(t^*) = B$ and
\[
\psi'(t^*) = -8B + 8v^\top A(t^*)^{2/3} v \le 8\bigl((dB)^{2/3} - B\bigr) < 0,
\]
since $B > d^2 \iff (dB)^{2/3} < B$. Hence $L(t^* - h) \ge \psi(t^* - h) > B$ for small $h > 0$, contradicting minimality of $t^*$. Letting $B \downarrow \max\{L(0), d^2\}$ gives $L(t) \le \max\{L(0), d^2\}$.

\emph{Forward completeness.} Taking square roots establishes~\eqref{eq:sv-envelope} on $[0, T_{\max})$. Let $\alpha_0 := \min\{\sigma_{\min}(W_0), c\} > 0$ and $\beta_0 := \max\{\sigma_{\max}(W_0), d\}$. Then the closed set
\[
\mathcal{K} := \{W \in \mathbb{R}^{m \times n} : \alpha_0 \le \sigma_{\min}(W) \le \sigma_{\max}(W) \le \beta_0\}
\]
is compact and contained inside $\mathcal{R}_{m,n}$. By standard continuation for real-analytic ODEs, $T_{\max} = \infty$.
\end{proof}

\section{Conjecture 2: Global Minima and Strict Saddles}

\begin{remark}[Typo in the printed stationary set of arXiv~v1]
In \href{https://arxiv.org/abs/2506.06134}{arXiv:2506.06134} (v1, \texttt{main.tex}, line~645), the stationary set $\mathcal{S}$ is printed with $(P V_C)_m^\top$ (permuting rows of $V_C$), whereas the paper's own derivation (\texttt{main.tex}, line~972) obtains $V_W = V_C P^\top$ (permuting columns).
For $m = 1, n = 2$, $V_C = \frac{1}{\sqrt{2}}\begin{psmallmatrix}1 & 1\\ 1 & -1\end{psmallmatrix}$, $D = \diag(2,1)$, $C = V_C D V_C^\top = \frac{1}{2}\begin{psmallmatrix}3 & 1\\ 1 & 3\end{psmallmatrix}$, and $P = \begin{psmallmatrix}0 & 1\\ 1 & 0\end{psmallmatrix}$, the printed formula gives $W = (PDP^\top)_{1,1}((PV_C)_{:,1})^\top = \frac{1}{\sqrt{2}}[1, 1]$, which satisfies $WW^\top = 1$, $WC = 2W$, and $-\nabla S_3(W) = 4(2^{2/3}-1)W \neq 0$ (norm $\approx 2.3496$).
Furthermore, if $V_C$ denotes one fixed eigenbasis rather than the union over all eigenbases of $C$, cutoff eigenvalue ties ($\lambda_m = \lambda_{m+1}$) produce continuous manifolds of global minima outside a single basis choice.
\end{remark}

\begin{theorem}[Landscape of $S_3$]\label{thm:c2}
Let $C \in \mathbb{R}^{n \times n}$ be symmetric positive definite and $1 \le m \le n$.
\begin{enumerate}[label=\textup{(\arabic*)}]
\item The stationary points of $S_3$ in $\mathcal{R}_{m,n}$ are precisely the matrices $W = U \Sigma V^\top$ with $U \in O_m$, $\Sigma = \diag(\sigma_1, \dots, \sigma_m) \succ 0$, and $V \in \mathbb{R}^{n \times m}$ satisfying $V^\top V = I_m$ and $CV = V\Sigma$. At any such point, $S_3(W) = -\sum_{i=1}^m \sigma_i^2$.
\item No rank-deficient matrix $W \in \mathbb{R}^{m \times n} \setminus \mathcal{R}_{m,n}$ is a local minimum of $S_3$. Every global minimum of $S_3$ on $\mathbb{R}^{m \times n}$ lies in $\mathcal{R}_{m,n}$, has value $-\sum_{i=1}^m \lambda_i^2$, and is characterized among full-rank stationary points by $\beta \le \alpha$, where $\alpha := \min_{1 \le i \le m} \sigma_i$ and $\beta := \lambda_{\max}(C|_{V^\perp})$ (taking $\beta = -\infty$ if $m = n$).
\item Every non-global full-rank stationary point ($\beta > \alpha$) is a strict saddle: its Hessian has a negative eigenvalue at most $4(1 - \beta/\alpha) < 0$ and a positive radial curvature $8/3 > 0$.
\end{enumerate}
\end{theorem}

\begin{proof}
\emph{Part~(1).} Set $K := (WCW^\top)^{1/3} \succ 0$. Stationarity $\nabla S_3(W) = 0$ is equivalent to $KW = WC$. Multiplying on the right by $W^\top$ gives $K(WW^\top) = WCW^\top = K^3$, hence $WW^\top = K^2$ and $K = (WW^\top)^{1/2}$. Writing a thin SVD $W = U \Sigma V^\top$ with $U \in O_m$, $\Sigma \succ 0$ diagonal, and $V^\top V = I_m$, we have $K = U\Sigma U^\top$, so $KW = WC \iff U\Sigma^2 V^\top = U\Sigma V^\top C \iff CV = V\Sigma$. Conversely, $CV = V\Sigma$ immediately implies $WCW^\top = U\Sigma^3 U^\top$, $K = U\Sigma U^\top$, $KW = WC$, and $S_3(W) = 2\tr(\Sigma^2) - 3\tr(\Sigma^2) = -\sum_{i=1}^m \sigma_i^2$.

\emph{Part~(2).} Suppose $\rank(W) = r < m$. Pick a unit vector $u \in \ker(W^\top) \subset \mathbb{R}^m$ ($u^\top W = 0$) and a unit vector $v \in \ker(WC) \subset \mathbb{R}^n$ (which exists because $\rank(WC) = r < m \le n$). Set $W_t := W + t u v^\top$ and $\gamma := v^\top C v > 0$. Because $u^\top W = 0$ and $WCv = 0$:
\[
\|W_t\|_F^2 = \|W\|_F^2 + 2t\,u^\top W v + t^2 = \|W\|_F^2 + t^2,
\qquad
W_t C W_t^\top = W C W^\top + t^2 \gamma\, u u^\top.
\]
Since $u \in \ker(W^\top) = \ker(WCW^\top)$, the rank-1 term $t^2 \gamma\, uu^\top$ is orthogonal to the range of $WCW^\top$, so its eigenvalue $t^2 \gamma$ appends cleanly to the nonzero spectrum of $WCW^\top$. Hence
\[
S_3(W_t) - S_3(W) = 2t^2 - 3|t|^{4/3}\gamma^{2/3} < 0
\qquad \text{for all } 0 < |t| < \Bigl(\frac{3}{2}\Bigr)^{3/2}\gamma,
\]
showing $W$ is not a local minimum. Since $S_3(W) \ge 2\|W\|_F^2 - 3\lambda_1^{2/3} m^{1/3} \|W\|_F^{4/3} \to +\infty$ as $\|W\|_F \to \infty$, $S_3$ attains its global minimum on $\mathbb{R}^{m \times n}$, necessarily at a full-rank stationary point $W = U\Sigma V^\top$. Because $C$ is symmetric and $CV = V\Sigma$, $V^\perp$ is $C$-invariant, so the eigenvalues of $C$ split into $\{\sigma_1, \dots, \sigma_m\}$ and the spectrum of $C|_{V^\perp}$. Thus $-\sum_{i=1}^m \sigma_i^2$ is minimal (equal to $-\sum_{i=1}^m \lambda_i^2$) if and only if every eigenvalue on $V^\perp$ is at most $\min_i \sigma_i$, i.e., $\beta \le \alpha$.

\emph{Part~(3).} Let $W = U\Sigma V^\top$ be a non-global full-rank stationary point, so $\beta > \alpha$. Choose $i$ with $\sigma_i = \alpha$, let $u_i$ be column $i$ of $U$, and choose a unit eigenvector $v_\perp \in V^\perp$ with $Cv_\perp = \beta v_\perp$. For $W_\varepsilon := W + \varepsilon u_i v_\perp^\top$, we have $W C v_\perp = \beta U \Sigma V^\top v_\perp = 0$ and $u_i^\top W = \alpha v_i^\top \perp v_\perp$, so $\|W_\varepsilon\|_F^2 = \|W\|_F^2 + \varepsilon^2$ and
\[
W_\varepsilon C W_\varepsilon^\top = U\diag(\sigma_1^3, \dots, \sigma_i^3 + \varepsilon^2 \beta, \dots, \sigma_m^3)U^\top.
\]
Hence
\[
S_3(W_\varepsilon) - S_3(W)
= 2\varepsilon^2 - 3\bigl((\alpha^3 + \varepsilon^2\beta)^{2/3} - \alpha^2\bigr)
= 2\Bigl(1 - \frac{\beta}{\alpha}\Bigr)\varepsilon^2 + O(\varepsilon^4),
\]
giving second directional derivative $\frac{d^2}{d\varepsilon^2}\big|_{\varepsilon=0} S_3(W_\varepsilon) = 4(1 - \beta/\alpha) < 0$.
Finally, along $W_\varepsilon = W + \varepsilon u_i v_i^\top$ (scaling $\sigma_i \mapsto \sigma_i + \varepsilon$), $S_3(W_\varepsilon) - S_3(W) = 2((\alpha+\varepsilon)^2 - \alpha^2) - 3(\alpha^{2/3}(\alpha+\varepsilon)^{4/3} - \alpha^2) = \frac{4}{3}\varepsilon^2 + O(\varepsilon^3)$, giving positive curvature $\frac{8}{3} > 0$. Thus $W$ is a strict saddle.
\end{proof}

\begin{remark}[Prior art on Theorem~\ref{thm:c2}]
The objective $S_3$ coincides with Eq.~(38) of Pehlevan, Sengupta, and Chklovskii (\href{https://arxiv.org/abs/1703.07914}{arXiv:1703.07914v2}, 2018), whose Theorem~1 and Eq.~(65) contain the corresponding $\tau \to 0$ instability block $2(\bar{M}^{-1}\sigma_{j+k} - I)$ matching $4(1-\beta/\alpha)$ after time rescaling. The corrected stationary set allowing arbitrary eigenvalue multiplicities also appears as Lemma~1 of Lipshutz and Lipshutz (\href{https://arxiv.org/abs/2601.13170}{arXiv:2601.13170}, 2026). The new contributions here are the rank-deficient exclusion proof, direct strict-saddle analysis on $S_3$ with arbitrary multiplicities, and identification of the arXiv~v1 typo.
\end{remark}

\section{Conjecture 3: Refutation and Exact Characterization}

Let $\Lambda_1 = [\diag(a_1, \dots, a_m),\ 0] \in \mathbb{R}^{m \times n}$ with $a_i \neq 0$ for all $i = 1, \dots, m$ (required for $\Lambda_1 U \Lambda_2 U^\top \Lambda_1^\top$ to be invertible so its $-1/3$ power exists), let $\Lambda_2 = \diag(d_1, \dots, d_n) \succ 0$ have distinct diagonal entries, and let $U \in O_n$. Conjecture~3 (\texttt{conj\_lin\_prog}) asserts that if
\[
H := U^\top \Lambda_1^\top \bigl(\Lambda_1 U \Lambda_2 U^\top \Lambda_1^\top\bigr)^{-1/3} \Lambda_1 U
\]
is diagonal, then $U$ is a permutation matrix.

\begin{theorem}[Exact diagonalization criterion and refutation]\label{thm:c3}
With $\Lambda_1, \Lambda_2, U$ as above:
\begin{enumerate}[label=\textup{(\arabic*)}]
\item $H$ is diagonal if and only if $K_U := U^\top (\Lambda_1^\top \Lambda_1) U$ is diagonal. In particular, diagonality of $H$ does not depend on $\Lambda_2$.
\item Whenever $n - m \ge 2$, taking $U = \diag(I_m, R)$ for any non-monomial $R \in O_{n-m}$ gives a strictly diagonal $H$ for \emph{every} $\Lambda_1$ and $\Lambda_2$.
\item Diagonality of $H$ forces $U$ to be a signed permutation matrix if and only if $|a_1|, \dots, |a_m|$ are pairwise distinct and $n - m \le 1$. (Sign flips always prevent forcing a pure permutation matrix.)
\end{enumerate}
\end{theorem}

\begin{proof}
Set $B := \Lambda_1 U \Lambda_2^{1/2} \in \mathbb{R}^{m \times n}$. Since every $a_i \neq 0$ and $\Lambda_2 \succ 0$, $B$ has full row rank $m$. Let $B = P S Q^\top$ be a thin SVD ($P \in O_m$, $S \in \mathbb{R}^{m \times m}$ positive diagonal, $Q \in \mathbb{R}^{n \times m}$ with $Q^\top Q = I_m$). Then $BB^\top = PS^2P^\top$ and $B^\top B = QS^2Q^\top$, whence
\[
\Lambda_2^{1/2} H \Lambda_2^{1/2}
= B^\top (BB^\top)^{-1/3} B
= Q S P^\top (P S^{-2/3} P^\top) P S Q^\top
= Q S^{4/3} Q^\top
= (B^\top B)^{2/3}.
\]
(Equivalently, with $K := (BB^\top)^{1/3}$ and $N := B^\top K^{-1} B$, one has $N^3 = (B^\top B)^2$.)
Because $X \mapsto X^{2/3}$ is a bijection on $\mathbb{S}^n_+$ preserving diagonal matrices (with inverse $Y \mapsto Y^{3/2}$) and $\Lambda_2^{1/2} \succ 0$ is diagonal,
\[
H \text{ diagonal}
\iff \Lambda_2^{1/2} H \Lambda_2^{1/2} \text{ diagonal}
\iff B^\top B \text{ diagonal}
\iff U^\top (\Lambda_1^\top \Lambda_1) U \text{ diagonal}.
\]
This proves~(1). Since $\Lambda_1^\top \Lambda_1 = \diag(a_1^2, \dots, a_m^2, 0, \dots, 0)$, any block-diagonal $U = \diag(I_m, R)$ with $R \in O_{n-m}$ commutes with $\Lambda_1^\top \Lambda_1$, proving~(2). Finally, $U^\top (\Lambda_1^\top \Lambda_1) U$ is diagonal if and only if the rows of $U$ permute and orthogonally mix only within eigenspaces of $\Lambda_1^\top \Lambda_1$. Every eigenspace is one-dimensional (forcing $U$ to be a signed permutation) if and only if $a_1^2, \dots, a_m^2$ are pairwise distinct and the zero eigenspace has dimension $n - m \le 1$, proving~(3).
\end{proof}

\begin{remark}[Survival of the paper's trace-maximization lemma]
Although Conjecture~3 fails across all five $(m, n)$ pairs tested in \href{https://arxiv.org/abs/2506.06134}{arXiv:2506.06134} (random orthogonal sampling misses block-diagonal kernel rotations because they have measure zero), the paper's trace-maximization lemma (\texttt{lem:pita}) still holds unconditionally: by Horn's log-majorization, the singular values of $\Lambda_1 U \Lambda_2^{1/2}$ are weakly log-majorized by the ordered products $|a|^\downarrow \odot (d^{1/2})^\downarrow$, and $x \mapsto \sum_{i=1}^m x_i^{4/3}$ is Schur-convex on $\mathbb{R}_{\ge 0}^m$. Thus the paper's alternative minimizer characterization survives by a different proof.
\end{remark}

\section{Formalization, Verification, and Attribution}

\paragraph{Lean~4 formalization.}
The file \texttt{SimMatchingPCA.lean} (in \texttt{lean\_proofs/Conjectures/}) compiles under Lean~\texttt{v4.24.0} with Mathlib, zero \texttt{sorry}s, and standard axioms only, machine-checking via an exact cube-root witness $K^3 = WCW^\top$:
\begin{itemize}[leftmargin=2em,itemsep=1pt,topsep=3pt]
\item Conjecture~1: \texttt{gram\_ode\_identity}, \texttt{lower\_barrier\_pos}, \texttt{upper\_barrier\_neg};
\item Conjecture~2: \texttt{stationary\_implies\_gram\_eq\_sq}, \texttt{printed\_set\_not\_stationary},\\
\texttt{rank\_one\_perturbation\_algebra}, \texttt{rank\_one\_escape\_negative},\\
\texttt{exchange\_curvature\_neg}, \texttt{radial\_curvature\_pos};
\item Conjecture~3: \texttt{conj3\_cube\_identity}, \texttt{conj3\_kernel\_rotation\_counterexample},\\
\texttt{conj3\_active\_tie\_counterexample}.
\end{itemize}
Full ODE global existence, Weyl/Dini continuation, general-dimension Hessian expansion, and the general Conjecture~3 iff are proved in prose above. Companion numerical stress checks live in \texttt{solutions/sim\_matching\_pca/explore.py}.

\paragraph{Literature note.}
All printed-statement checks refer to arXiv~v1 (June~2025); the published \emph{Neural Computation} 38(5):725--764 (April~2026) pages returned HTTP~403 and were not directly inspected. Related landscape results for Conjecture~2 appear in Pehlevan--Sengupta--Chklovskii (\href{https://arxiv.org/abs/1703.07914}{arXiv:1703.07914v2}) and Lipshutz--Lipshutz (\href{https://arxiv.org/abs/2601.13170}{arXiv:2601.13170}); no prior resolution of Conjecture~1 or refutation of Conjecture~3 was found as of November~2026.

\end{document}
